\bookinterlude{II}{III}{A pergunta a lápis}
Encomendamos uma bucha à oficina do outro lado do pátio. A ficha volta com uma pergunta a lápis. “De qual torno vai sair a peça?” Até ontem eu teria achado graça. Só tínhamos um.

No nosso torno, as marcas 1 e 2 do volante deram peças de 5 e 6 mm. No deles, deram 9 e 8 mm, nessa ordem. Junto as fichas e ajusto por mínimos quadrados uma reta para prever o diâmetro pela posição do volante. Ela prevê 7 mm nas duas marcas. Ninguém fez uma peça de 7 mm.

Tiro uma das quatro fichas ao acaso. Pelo telefone, você quer saber se a peça cabe num furo de 7 mm. “Estava na marca 1?” Um sim deixa você entre 5 e 9 mm. Um não, entre 6 e 8. Comece pela pergunta a lápis e uma resposta já lhe diz se cabe. Se quiser também o diâmetro, guarde a pergunta sobre o volante para depois.

Amanhã chegam os cadernos de oito oficinas, com colunas para o material, o turno e a última troca de ferramenta. Quando o telefone tocar, como vai escolher, pelas fichas de ontem, a primeira pergunta de hoje?
